\scanpage{227}
\mspage{143}
% archive sheet number: 226
% element: text | top, corollary conclusion | alors on a
alors on a $\operatorname{card\,réd}C\leq\pi_1^{\pi_0}$.
% element: margin note | upper left, slanted | majoration non prouvee
% slanted
\margin{majoration \unsure{non prouvée}, \unsure{prob.blemt|prob.blement|probablement} fausse, cf \unsure{annotation} ci-dessus.}

% element: text | upper third, proof | Le 3 se prouve
\unsure{Le|La} 3°) se prouve par la démonstration \deleted{\unsure{préc.|puis}} transfinie précédente, \deleted{\unsure{on|les}} on prouve par récurrence transfinie que
% element: display | upper third, inductive bound | card red C alpha
\[
\operatorname{card\,réd}C_\alpha\leq\pi_1^{\pi_0}
\]
% element: text | upper middle, initial and successor cases | pour tt alpha
pour \unsure{tt|$\forall$} $\alpha$. C'est \unsure{o.K|o.k} pour $\alpha=0$ par hyp (puisque $\pi_1\leq\pi_1^{\pi_0}$), \deleted{c'est vrai} si c'est vrai pour $C_\alpha$, alors 2°) assure que
% element: display | middle, successor estimate | card red C alpha power estimate
\[
\operatorname{card\,red}C_\alpha\leq(\pi_1^{\pi_0})^{\pi_0}=\pi_1^{\pi_0\cdot\pi_0}=\pi_1^{\pi_0},
\]
% element: text | middle, limit case | enfin si alpha est un ordinal limite
enfin, si $\alpha$ est un ordinal limite $<\alpha_0$, on a $C_\alpha=\bigcup_{\alpha'<\alpha}C_{\alpha'}$, donc
% element: display | middle lower, limit estimate | card red C alpha product estimate
\[
\operatorname{card\,réd}C_\alpha\leq\pi_0\cdot\pi_1^{\pi_0}=\pi_1^{\pi_0}
\]
% element: text | lower middle, justification and successor cardinal | vu que
(\unsure{vu|car} que $\pi_1^{\pi_0}\geq\pi_0$), \deleted{Donc \unsure{on a|il est}} \inserted{en vertu de} \deleted{par} 1° ci-dessus, Et c'est dans ce cas que $\operatorname{card}I_\alpha\leq\pi_0$. Soit $\pi'_0$ le cardinal \inserted{successeur} de $\pi_0$, alors \unsure{la|le} $\unsure{=|\bar m}$ majoration nous donne
% element: display | lower third, bound at alpha zero | card red C alpha zero
\[
\operatorname{card\,red}C_{\alpha_0}\leq\pi'_0\pi_1^{\pi_0}.
\]
% element: text | lower third, cardinal comparison | Mais
Mais \unsure{on a|...} $\pi_1^{\pi_0}\geq 2^{\pi_0}>\pi_0$, donc $\pi_1^{\pi_0}\geq\pi'_0$, donc $\pi'_0\pi_1^{\pi_0}=\pi_1^{\pi_0}$, d'où
% element: display | foot, final estimate | card red C alpha zero final bound
\[
\operatorname{card\,red}C_{\alpha_0}\leq\pi_1^{\pi_0},
\]
% element: text | foot, conclusion | ce qui est
ce qui est la majoration annoncée.
